\chapter{Introduction}
\label{ch:introduction}

\section{Motivation and Scope}
\label{sec:motivation_scope}

Organizations are expected to deliver essential services despite operational disruptions that may affect their activities, infrastructure, information, or external dependencies. The consequences of an interruption depend not only on the nature of the incident but also on the importance of the affected service and the duration for which it remains unavailable. An interruption that is manageable initially may become unacceptable as its effects on service users, operations, financial obligations, regulatory requirements, or institutional reputation accumulate. Consequently, effective business continuity management requires organizations to identify their priority services, understand how interruption impacts develop over time, and establish defensible recovery requirements before a disruption occurs \cite{nrcnationalframework,dga2025bcm}.

Business continuity management (BCM) provides a structured approach to preparing for disruptions and maintaining or restoring prioritized activities and services. Business impact analysis (BIA) is central to this approach because it examines the consequences of interruption, establishes the relative importance of services and activities, and informs recovery planning. Two recovery parameters are particularly relevant to this research: the maximum tolerable period of disruption (MTPD), which expresses the duration beyond which the consequences of an interruption become unacceptable, and the recovery time objective (RTO), which specifies a target time for resuming a service or activity. Although related, these parameters serve different purposes and must be established with reference to the organization's assessed impacts and recovery requirements \cite{nrcnationalframework,dga2025bcm}.

The assessment of interruption consequences is inherently time-dependent. Services differ in the speed at which interruption impacts become serious, and the criteria used to judge those impacts may also differ among organizations. The Digital Government Authority's business continuity guidance, for example, presents impact-based classifications and recovery targets for government digital platforms, while recognizing that assessment time intervals should reflect the sensitivity of the activities concerned \cite{dga2025bcm}. These principles motivate an approach in which an organization can define its own assessment periods, impact criteria, acceptable thresholds, and corresponding recovery rules instead of assuming that a single fixed configuration is universally appropriate.

This project proposes \emph{Time-Dependent Recovery Prioritization} (TDRP), a configurable rule-based methodology, and \emph{Masar}, a web-based decision support system (DSS) that implements the methodology. The proposed approach begins with business services and their BIA results. It applies explicitly defined rules to time-dependent impact assessments to determine service criticality and recovery requirements and then uses these results to establish relative service recovery priorities. Once services have been assessed, their supporting dependencies can be identified and related to the criticality of the services they enable. The system records the relevant relationships and presents the resulting assessments and priorities through an operational dashboard. Its purpose is to support the interpretation and use of continuity information by responsible personnel, not to replace professional judgment or assume responsibility for operational recovery decisions.

The scope of the research is the design, implementation, and evaluation of the TDRP methodology within Masar. Its central concern is the service-level assessment and prioritization process; supporting resources are considered insofar as their relationships help explain service delivery and recovery dependencies. 

\section{Problem Statement and Significance for Information Systems}
\label{sec:problem_significance}

Determining recovery priorities requires more than identifying that a service is important. Organizations must explain how a service's assessed interruption impacts justify its criticality, how those impacts relate to acceptable disruption durations and recovery targets, and how the resulting priorities should be interpreted alongside the dependencies required for service delivery. National business continuity guidance identifies BIA, recovery requirements, and dependencies as important components of continuity management \cite{nrcnationalframework,dga2025bcm}. Nevertheless, these principles do not establish a single mandatory computational procedure that all organizations must use to translate successive impact assessments into MTPD, RTO, and service recovery priorities. Nor should one organization's specific thresholds or assessment intervals be presumed suitable for every other organization.

The research problem is therefore the need for a clear, configurable, and reproducible method for translating time-dependent business-service impact assessments into recovery requirements and relative recovery priorities, while retaining an explicit connection between assessed services and the dependencies that support their delivery. A useful method must make its decision rules visible, permit different organizations to configure those rules according to their approved criteria, and produce results that can be traced to the underlying assessments. It must also distinguish the priority assigned to a service from the importance attributed to a supporting dependency because of the services that rely on it. The project does not assume that existing continuity practices are absent or that organizations lack information; rather, it addresses the formalization and implementation of a particular service-first prioritization procedure.

This problem is significant to information systems because it concerns the transformation of organizational assessments into consistent and explainable decision-support outputs. A DSS can represent services, time-based impacts, recovery parameters, configurable rules, and dependency relationships within a coherent information model, then apply the same declared logic to comparable inputs. This makes it possible to inspect how a priority was obtained, examine the effect of changing approved rule configurations, and present the resulting information to continuity practitioners. The design and evaluation of such a purposeful information-systems artifact are consistent with design-science research, which focuses on constructing and assessing artifacts that address identified organizational problems \cite{hevner2004design,peffers2007dsrm}.

\section{Importance of the Problem and the Big Picture}
\label{sec:importance}

Business continuity is an organizational capability rather than a purely technical recovery activity. A disruption can affect the availability of public or private services, the obligations of an organization, the experience of beneficiaries, and its ability to maintain essential operations. The Saudi National Framework places business continuity within a wider context of risk management, emergency management, preparedness, response, and organizational resilience, while the Digital Government Authority emphasizes sustained service delivery and the alignment of recovery arrangements with organizational objectives \cite{nrcnationalframework,dga2025bcm}. From this perspective, prioritization is important because organizations must determine which service outcomes are most time-sensitive and how recovery efforts should be informed by those priorities.

A service-oriented perspective also provides a principled way to consider supporting dependencies. Services are delivered through interacting activities and enabling resources, and interruption of a shared dependency may affect more than one service. The criticality of a supporting element should therefore be interpreted in relation to the services it enables, rather than assumed from its characteristics alone. In the proposed methodology, service criticality is established first, after which relevant dependency relationships are recorded and a stated rule is used to associate supporting elements with service-level criticality. Where several services share a dependency, the methodology must define how their classifications are combined. Using the highest criticality among supported services is a proposed rule to be examined and validated, not a universal requirement attributed to the reference frameworks.

A further concern is the balance between standardization and organizational flexibility. Consistency requires a common method for applying rules, whereas practical use requires organizations to select impact scales, assessment time frames, and recovery thresholds that reflect their contexts. A configurable rule-based methodology addresses both concerns by standardizing the form and execution of the decision logic without imposing identical parameter values on every organization. In this sense, the anticipated contribution is not to replace existing BCM standards or practitioners' expertise, but to provide a transparent procedure for interpreting BIA results and expressing recovery priorities. Research on business impact analysis and integrated continuity planning provides a broader basis for investigating this connection between organizational impacts and recovery planning \cite{torabi2014framework,sahebjamnia2015integrated}.

\section{Objectives and Proposed Solution}
\label{sec:objectives_solution}

The aim of this project is to design, implement, and evaluate a configurable rule-based methodology for deriving recovery requirements and prioritizing critical business services using time-dependent impact assessments. The methodology, TDRP, is the main research contribution. Masar is the proposed web-based DSS through which the methodology will be configured, applied, and evaluated, and through which its results will be presented to authorized users.

The methodology will represent BIA findings at the service level across successive assessment periods and relevant impact dimensions. Based on the organization's defined criteria and decision rules, it will determine service criticality, identify the applicable MTPD, and establish the corresponding RTO. The method will document how temporal assessments, impact thresholds, and recovery-time mappings are interpreted, including the handling of incomplete data and cases in which an assessed impact threshold is not reached. Organizations will be able to configure the permitted parameters to reflect their approved continuity requirements.

Using these derived service-level results, TDRP will establish relative recovery priorities through explicit rules that consider criticality and recovery-time requirements. The approach will then associate assessed services with the dependencies necessary for their delivery. A supporting element will inherit criticality from the service or services it supports according to a documented rule; for a shared element, the proposed default is the highest criticality among its associated services, subject to evaluation. This association is intended to preserve the business rationale for the importance of supporting elements. It does not equate a service-priority ranking with an automatically executable sequence of recovery operations.

Masar will implement the methodology through an integrated data model and rule-processing functions. It will enable authorized users to maintain services, impact assessments, configuration settings, recovery requirements, and relevant dependency relationships. An operational dashboard will present the resulting service classifications, MTPD and RTO values, recovery priorities, and linked supporting elements. During an actual interruption, users may consult the recorded information for affected services to understand their established priority and dependencies. The dashboard will display the methodology's outputs and associated context; it will not simulate incidents, infer operational actions, autonomously recommend recovery procedures, or make or execute recovery decisions.

The research will evaluate both the methodological rules and their implementation. Evaluation will examine the consistency and reproducibility of the outputs, the behavior of the method under different valid configurations, the correctness of derived MTPD and RTO values against specified test cases, the treatment of shared dependencies, and the traceability of displayed priorities to the originating impact assessments and rules. The intended deliverables are a documented TDRP methodology, an implemented Masar prototype, and an evaluation of their correctness and usefulness as a basis for structured decision support. Operational effectiveness in real incidents remains outside the claims that can be established by prototype-level evaluation alone.
